% Vietnamese translation for OI Public Library
% Source-File: IOI中国国家候选队论文/2003/何林/何林.doc
\documentclass[11pt]{article}
\usepackage{oipl-vn}

\title{Một cách giải cho lớp bài toán cân bi}
\author{He Lin}
\date{2003}

\begin{document}
\maketitle

\origtitle{A solution for a class of ball-weighing problems}
\sourcepath{IOI China candidate team papers / 2003 / He Lin / He Lin.doc}

\begin{abstract}
Bài viết đưa ra một cách giải đầy đủ và chặt chẽ cho một lớp bài toán cân bi
bằng cân hai đĩa. Từ đó, tác giả tổng kết một số kinh nghiệm trong quá trình
nghiên cứu: bắt đầu từ bài toán đơn giản, dùng cây quyết định để chứng minh
cận dưới, rồi xây dựng phương án đạt cận bằng cách phân bố đều các khả năng.
\end{abstract}

\noindent\textbf{Từ khóa:} cây quyết định, chia ba, phân bố đều.

\section{Dẫn nhập}

Có \(n\) viên bi, trong đó đúng một viên là bi lỗi. Ta có một chiếc cân hai
đĩa và muốn xác định viên bi lỗi. Xét bốn biến thể sau.

\begin{enumerate}
  \item Bi lỗi nặng hơn các viên còn lại.
  \item Không biết bi lỗi nặng hơn hay nhẹ hơn, nhưng chỉ cần tìm ra viên bi
  lỗi.
  \item Không biết bi lỗi nặng hơn hay nhẹ hơn; cần tìm ra viên bi lỗi và xác
  định nó nặng hay nhẹ.
  \item Như bài toán 3, nhưng ta còn có thêm một viên bi chuẩn.
\end{enumerate}

Những bài toán này có hình thức rất gần nhau. Câu hỏi chính là nên bắt đầu
phân tích từ đâu, và kinh nghiệm ở bài toán đơn giản có thể được chuyển sang
bài toán phức tạp đến mức nào.

\section{Bài toán đơn giản}

Trước hết xét bài toán 1. Đánh số các viên bi theo thứ tự
\(1,2,\ldots,n\). Với \(n=3\), đặt viên \(1\) lên đĩa trái và viên \(2\) lên
đĩa phải. Nếu cân nghiêng trái thì viên \(1\) là bi lỗi; nếu nghiêng phải thì
viên \(2\) là bi lỗi; nếu cân bằng thì hai viên \(1,2\) đều chuẩn, nên viên
\(3\) là bi lỗi. Vì vậy \(f(3)=1\).

Với \(n=9\), chia các viên thành ba nhóm
\[
  A=\{1,2,3\},\qquad B=\{4,5,6\},\qquad C=\{7,8,9\}.
\]
Cân \(A\) với \(B\). Nếu cân nghiêng trái, bi lỗi nằm trong \(A\); nếu nghiêng
phải, bi lỗi nằm trong \(B\); nếu cân bằng, bi lỗi nằm trong \(C\). Trong mọi
trường hợp, phạm vi nghi vấn giảm từ \(9\) xuống \(3\), nên \(f(9)=2\).

\paragraph{Định lý 1.}
Với \(k\ge 1\), \(f(3^k)=k\).

\paragraph{Chứng minh.}
Các trường hợp nhỏ đã xét. Giả sử biết cách giải với \(3^k\) viên trong
\(k\) lần cân. Với \(3^{k+1}\) viên, chia đều thành ba nhóm \(A,B,C\), mỗi
nhóm \(3^k\) viên, rồi cân \(A\) với \(B\). Sau một lần cân, bi lỗi chỉ còn
có thể nằm trong đúng một nhóm. Theo giả thiết quy nạp, cần thêm \(k\) lần cân
nữa. Tổng cộng là \(k+1\) lần.

Suy ra với \(n\) bất kỳ,
\[
  f(n)=\left\lceil \log_3 n \right\rceil .
\]
Cách làm là mỗi lần chia tập nghi vấn thành ba phần càng đều càng tốt.

Để chứng minh tối ưu, ta dùng \term{cây quyết định}. Mỗi lá biểu diễn một kết
quả cuối cùng có thể xảy ra; mỗi nút trong biểu diễn một lần cân; mỗi nút
trong có nhiều nhất ba con, tương ứng với ba kết quả: nghiêng trái, nghiêng
phải, cân bằng. Mọi phương án cân đều có thể biểu diễn bằng một cây quyết
định, và độ sâu của cây chính là số lần cân trong trường hợp xấu nhất.

Trước khi cân, mỗi viên đều có thể là bi lỗi, nên cây quyết định có ít nhất
\(n\) lá. Một cây ba nhánh độ sâu \(h\) có nhiều nhất \(3^h\) lá, do đó
\[
  h\ge \left\lceil \log_3 n \right\rceil .
\]
Cận dưới này trùng với phương án chia ba ở trên, vì vậy bài toán 1 đã được
giải trọn vẹn.

Điểm đáng chú ý là ``chia ba'' không phải một mẹo cảm tính. Nó đến từ cấu
trúc ba nhánh của cây quyết định. Đồng thời, chữ ``đều'' cũng rất quan trọng:
độ sâu của một cây phụ thuộc vào cây con sâu nhất của gốc, nên muốn tối ưu
trong trường hợp xấu nhất thì số lá phải được phân bố càng đều càng tốt.

\section{Mở rộng: có bi chuẩn và cần biết nặng nhẹ}

Xét bài toán 4. Có \(n\ge 3\) viên, đúng một viên là bi lỗi; bi lỗi có trọng
lượng khác các viên chuẩn nhưng chưa biết nặng hơn hay nhẹ hơn. Ta có một cân
hai đĩa và một viên bi chuẩn. Cần tìm bi lỗi và xác định nó nặng hay nhẹ.

Nếu chỉ bắt chước bài toán 1, ta sẽ gặp ngay trở ngại: khi cân hai nhóm với
nhau và cân lệch, bi lỗi có thể là một viên nặng ở bên nặng hơn, hoặc một
viên nhẹ ở bên nhẹ hơn. Sau một lần cân, cấu trúc của tập nghi vấn không còn
là ``một tập các viên tương đương'' nữa. Ta cần một mệnh đề mạnh hơn.

\paragraph{Bổ đề.}
Cho hai nhóm bi. Nhóm thứ nhất có \(n\) viên, mỗi viên nếu lỗi thì chỉ có thể
là nặng; nhóm thứ hai có \(m\) viên, mỗi viên nếu lỗi thì chỉ có thể là nhẹ.
Biết rằng đúng một viên trong hai nhóm là bi lỗi. Khi \(n+m\ge 2\), có thể
tìm viên lỗi trong
\[
  \left\lceil \log_3(n+m) \right\rceil
\]
lần cân, và đây là tối ưu.

\paragraph{Chứng minh bổ đề.}
Cây quyết định phải có \(n+m\) lá, nên cận dưới là
\(\lceil\log_3(n+m)\rceil\).

Ta chứng minh có thể đạt cận này bằng quy nạp theo \(n+m\). Ký hiệu \(S\) là
bi chuẩn; \(iH\) nghĩa là viên \(i\) nặng; \(iL\) nghĩa là viên \(i\) nhẹ.
Ký hiệu \(\langle A,B\rangle\) là kết quả khi đặt tập \(A\) lên đĩa trái và
tập \(B\) lên đĩa phải; kết quả có thể là \(\mathrm{Left}\),
\(\mathrm{Right}\) hoặc \(\mathrm{Middle}\).

Trường hợp \(n+m=2\) chỉ cần một lần cân. Với bước quy nạp, giả sử
\(n+m=k\). Ta chia các khả năng thành ba phần đều nhất có thể. Ví dụ, nếu
\(n=3p\) và \(m=3q\), chia nhóm nặng-khả-dĩ thành \(A,B,C\), mỗi nhóm \(p\)
viên; chia nhóm nhẹ-khả-dĩ thành \(A',B',C'\), mỗi nhóm \(q\) viên. Cân
\(\langle A\cup A', B\cup B'\rangle\).

\begin{itemize}
  \item Nếu cân bằng, bi lỗi nằm trong \(C\cup C'\).
  \item Nếu nghiêng trái, bi lỗi nằm trong \(A\cup B'\): hoặc một viên nặng
  trong \(A\), hoặc một viên nhẹ trong \(B'\).
  \item Nếu nghiêng phải, bi lỗi nằm trong \(A'\cup B\).
\end{itemize}

Mỗi nhánh còn đúng \(p+q\) khả năng, nên theo giả thiết quy nạp cần thêm
\(\lceil\log_3(p+q)\rceil\) lần cân.

Các trường hợp số dư khác của \(n,m\) theo modulo \(3\) được xử lý bằng cùng
nguyên tắc: không nhất thiết chia đều số viên, mà chia đều số khả năng. Một
bảng chia điển hình là
\[
\begin{array}{c|c|c}
\text{trường hợp} & (|A|,|B|,|C|) & (|A'|,|B'|,|C'|)\\ \hline
n=3p,\ m=3q+1 & (p,p,p) & (q,q,q+1)\\
n=3p,\ m=3q-1 & (p,p,p) & (q,q,q-1)\\
n=3p+1,\ m=3q+1 & (p,p+1,p) & (q+1,q,q)\\
n=3p+1,\ m=3q+2 & (p,p,p+1) & (q+1,q+1,q)
\end{array}
\]
Mục tiêu của bảng là bảo đảm ba kết quả cân nhận số lá gần bằng nhau nhất.
Nhờ vậy nhánh xấu nhất cũng đạt cận tối ưu.

Quay lại bài toán 4. Vì mỗi viên có thể là nặng hoặc nhẹ, có \(2n\) khả năng
cuối cùng. Do đó bất kỳ cây quyết định nào cũng có độ sâu ít nhất
\[
  \left\lceil \log_3(2n) \right\rceil .
\]
Ta sẽ xây dựng phương án đạt cận này.

Với \(n=2\), có thể kiểm tra trực tiếp bằng một cây quyết định độ sâu \(2\).
Giả sử kết luận đúng với mọi số viên nhỏ hơn \(k\).

Nếu \(k=3p\), chia thành ba nhóm \(A,B,C\), mỗi nhóm \(p\) viên, rồi cân
\(\langle A,B\rangle\). Nếu cân lệch, bài toán trở thành bổ đề với \(p+p\)
khả năng; nếu cân bằng, bi lỗi nằm trong \(C\) và ta dùng giả thiết quy nạp.

Nếu \(k=3p+2\), chọn \(|A|=|B|=p+1\), \(|C|=p\). Ba nhánh có lần lượt
\(2p+2,2p+2,2p\) khả năng, đủ đều để đạt cận tối ưu.

Trường hợp dễ sai nhất là \(k=3p+1\). Nếu chia thành \(p,p,p+1\) viên thì ba
nhánh có \(2p,2p,2p+2\) khả năng, không đều. Cách đúng là đưa bi chuẩn vào
một đĩa:
\[
  A=\{S,1,2,\ldots,p\},\quad
  B=\{p+1,p+2,\ldots,2p+1\},\quad
  C=\{2p+2,\ldots,3p+1\}.
\]
Khi cân \(\langle A,B\rangle\), bi chuẩn không phải ứng viên lỗi, nên nếu cân
lệch thì còn \(2p+1\) khả năng; nếu cân bằng thì còn \(2p\) khả năng. Đây là
cách phân bố lá đúng theo cây quyết định.

Vì vậy, với một bi chuẩn, số lần cân cần và đủ để tìm bi lỗi đồng thời xác
định nặng nhẹ là
\[
  \left\lceil \log_3(2n) \right\rceil .
\]

\section{Không có bi chuẩn}

Xét bài toán 3, tức là cũng cần biết bi lỗi nặng hay nhẹ nhưng không có sẵn
bi chuẩn. Sau lần cân đầu tiên, ta thường có thể thu được bi chuẩn: nếu cân
bằng thì các viên đã cân đều chuẩn; nếu cân lệch thì các viên không được cân
đều chuẩn. Vì vậy từ lần thứ hai trở đi, bài toán gần như trở về bài toán 4.

Khác biệt nằm ở lần cân đầu tiên. Khi \(n\bmod 3=0\) hoặc \(n\bmod 3=2\), ta
không cần dùng bi chuẩn để đạt phân bố tối ưu. Riêng \(n\bmod 3=1\), nếu
không có bi chuẩn thì lần cân đầu tiên không thể phân bố hoàn hảo các lá của
cây quyết định; ta phải dùng một phương án kém hơn một nhánh.

Kết luận cổ điển là: không có bi chuẩn, để tìm bi lỗi và biết nó nặng hay nhẹ
cần số lần cân nhỏ nhất là
\[
  \left\lceil \log_3(2n+3) \right\rceil .
\]
Sự xuất hiện của \(+3\) phản ánh đúng mất mát ở lần cân đầu tiên khi thiếu
bi chuẩn.

\section{Chỉ cần tìm viên lỗi}

Bây giờ xét bài toán 5: có một bi chuẩn, bi lỗi có thể nặng hoặc nhẹ, nhưng
ta chỉ cần biết viên nào là lỗi, không cần biết nó nặng hay nhẹ.

Lúc này cận dưới không còn hiển nhiên. Về mặt trạng thái vật lý có \(2n\) khả
năng, nhưng kết quả cuối cùng chỉ cần phân biệt \(n\) viên. Ta dùng một mô
hình đệ quy.

Gọi \(f(n)\) là số lần cân ít nhất. Giả sử tạm thời có nhiều bi chuẩn. Ở lần
cân đầu, lấy \(a\) viên nghi vấn đặt bên trái và \(b\) viên nghi vấn đặt bên
phải, với \(a\le b\); bù thêm \(b-a\) bi chuẩn vào bên trái để hai đĩa có
cùng số viên.

Nếu cân lệch, bi lỗi nằm trong \(a\) viên bên trái theo khả năng nặng, hoặc
trong \(b\) viên bên phải theo khả năng nhẹ. Theo bổ đề, cần thêm
\(\lceil\log_3(a+b)\rceil\) lần cân. Nếu cân bằng, bi lỗi nằm trong
\(n-a-b\) viên còn lại, cần thêm \(f(n-a-b)\) lần cân. Do đó
\[
  f(n)=\min_{1\le p\le n}
  \left(
    \max\left\{\left\lceil\log_3 p\right\rceil,\ f(n-p)\right\}+1
  \right),
\]
trong đó \(p=a+b\) và có thể chọn \(a,b\) sao cho \(|a-b|\le 1\), nên thật ra
chỉ cần nhiều nhất một bi chuẩn.

Từ các giá trị ban đầu
\[
  f(0)=0,\qquad f(1)=0,\qquad f(2)=1,
\]
ta thu được bảng:
\[
\begin{array}{c|c}
n & f(n)\\ \hline
1 & 0\\
2 & 1\\
3\text{--}5 & 2\\
6\text{--}14 & 3\\
15\text{--}41 & 4\\
42\text{--}122 & 5
\end{array}
\]
Mẫu số liệu gợi ý công thức
\[
  f(n)=\left\lceil \log_3(2n-1) \right\rceil .
\]

Có thể chứng minh công thức này bằng quy nạp. Cận dưới lấy trực tiếp từ truy
hồi: nếu tồn tại \(p\) làm cho \(f(k)<\lceil\log_3(2k-1)\rceil\), hai điều
kiện ở nhánh lệch và nhánh cân bằng sẽ dẫn đến hai bất đẳng thức mâu thuẫn
khi cộng lại. Cận trên được xây dựng theo ba trường hợp của \(k\bmod 3\):

\begin{itemize}
  \item Nếu \(k=3p\), chia thành ba nhóm \(p,p,p\) và cân hai nhóm đầu.
  \item Nếu \(k=3p+1\), dùng bi chuẩn trong nhóm thứ nhất để tạo hai đĩa kích
  thước \(p+1\), còn nhóm ngoài cân có \(p\) viên nghi vấn.
  \item Nếu \(k=3p+2\), cũng dùng bi chuẩn để cân hai nhóm kích thước
  \(p+1\), còn lại \(p+1\) viên.
\end{itemize}

Trong mỗi trường hợp, nhánh cân bằng trở về bài toán 5 với ít viên hơn, còn
nhánh lệch trở về bổ đề hai nhóm. Số khả năng ở các nhánh được phân bố đều
theo đúng công thức trên. Vì vậy, với một bi chuẩn và chỉ cần tìm viên lỗi,
số lần cân tối ưu là
\[
  \left\lceil \log_3(2n-1) \right\rceil .
\]

Bài toán 2 là biến thể bỏ bi chuẩn khỏi bài toán 5. Tác giả gợi ý người đọc
tự suy luận bằng cùng kỹ thuật.

\section{Tổng kết phương pháp}

Với \(n\ge 3\), các kết luận chính là:

\begin{enumerate}
  \item Nếu biết bi lỗi nặng hơn, cần
  \(\lceil\log_3 n\rceil\) lần cân để tìm bi lỗi.
  \item Nếu không biết nặng nhẹ, có một bi chuẩn, và cần biết cả nặng nhẹ, cần
  \(\lceil\log_3(2n)\rceil\) lần cân.
  \item Nếu không biết nặng nhẹ, không có bi chuẩn, và cần biết cả nặng nhẹ,
  cần \(\lceil\log_3(2n+3)\rceil\) lần cân.
  \item Nếu không biết nặng nhẹ, có một bi chuẩn, nhưng chỉ cần tìm viên lỗi,
  cần \(\lceil\log_3(2n-1)\rceil\) lần cân.
\end{enumerate}

Tư tưởng thống nhất là dùng cây quyết định làm cầu nối giữa phương án cân và
chứng minh tối ưu. Khi cấu trúc cây đã rõ, nguyên tắc cốt lõi là phân bố đều
các lá, tức là phân bố đều các khả năng còn lại, chứ không nhất thiết là chia
đều số viên bi.

Từ quá trình nghiên cứu này có thể rút ra vài kinh nghiệm:

\begin{enumerate}
  \item Bắt đầu từ trường hợp đơn giản, rồi từng bước thêm điều kiện khó hơn.
  \item Tổng kết kinh nghiệm từ bài toán đơn giản và ngoại suy một cách hợp lý.
  \item Khi điều kiện thay đổi, phải xem xét phần khác biệt của cấu trúc bài
  toán, không chỉ nhìn vào hình thức giống nhau.
  \item Có thể mạnh dạn phỏng đoán, nhưng phải dùng chứng minh chặt chẽ để
  xác nhận.
\end{enumerate}

\section*{Phụ lục và tài liệu tham khảo}

Tác giả có kèm một chương trình mô phỏng đơn giản cho bài toán 3 trong
\texttt{Game.pas}; bản dịch này không dịch mã chương trình.

Cây quyết định còn có thể chứng minh rằng mọi thuật toán sắp xếp dựa trên so
sánh đều cần ít nhất
\[
  \left\lceil \log_2(n!) \right\rceil
\]
phép so sánh trong trường hợp xấu nhất. Theo công thức Stirling,
\(\log_2(n!)=O(n\log n)\), nên \(O(n\log n)\) là bậc cận dưới lý thuyết cho
sắp xếp bằng so sánh.

\begin{thebibliography}{9}
\bibitem{yan}
Yan Weimin, Wu Weimin.
\textit{Data Structures}, second edition.
Tsinghua University Press.
\end{thebibliography}

\end{document}
